\section{Method}\label{sec:method}

\subsection{Twist from a tracked
path}\label{sec:twist}

A 3 s segment at constant drive is, after the lag has settled, a
circular arc: in the body frame the marker moves with constant twist
\((u_1,u_2,\omega)\), and the instantaneous centre of rotation sits at
\((-u_2/\omega,\ u_1/\omega)\). We rotate the differenced marker track
into the body frame, where the velocity is linear in \((u_1, u_2)\) and
the heading increment in \(\omega\), and fit the twist by linear least
squares over the steady part of the segment; on the kinematic model of
the lab's twin this recovers \(u_1\) to within 0.3 mm/s in simulation.
The first 0.3 s after the command and the stop transient are excluded;
they are used in §\ref{sec:lag} for the lag.

\subsection{Forward map}\label{sec:forward-map}

Each motor's drive \(p\in\{L,R\}\) enters through an activation with a
dead zone and a power law,
\[ f_k = \max(0, |p_k| - p_{0,k})^{\gamma_k}, \qquad k \in \{L,R\}, \]
and the twist is linear in the activations:
\[ \mathbf y = \mathbf C\,\boldsymbol\phi(f_L, f_R, \Delta V), \qquad \boldsymbol\phi_\mathrm{lin} = (1, f_L, f_R)^\top, \]
with variants \texttt{int} (adds \(f_L f_R\)), \texttt{bat} (adds
\(f_L\Delta V, f_R\Delta V\) with \(\Delta V\) the battery voltage
relative to a reference) and \texttt{cor} (§\ref{sec:identifiability}).
The activation uses \(|p_k|\), so the map covers the positive drive
quadrant; the eight reverse-polarity cells of the protocol are checks of
the firmware's motor directions, not part of the fit. A segment in which
both motors are in their dead zone predicts zero twist. The dead zones
and exponents are shared by the three outputs because they describe the
motors, not the body; given them, \(\mathbf C\) is a linear
least-squares solution, so the fit is a Nelder--Mead search over four
nonlinear parameters with the coefficients projected out (variable
projection), with residuals scaled per output so that 0.5 cm/s and 0.1
rad/s count alike. Variants, and a Gaussian process and a small neural
network as non-parametric comparators, are compared by grouped
cross-validation in which all replicates of a drive cell are held out
together, so that the score measures interpolation to unseen drive
levels rather than replicate noise.

\subsection{Identifiability of the centre-of-rotation
offsets}\label{sec:identifiability}

The twin's kinematic model, and GRASPion's ``rotation about a front
leg'', place the centre of rotation at a lateral offset \(d\) from the
reference point, with separate values \(d_\mathrm{cw}\) and
\(d_\mathrm{ccw}\) for the two turning directions, so that the forward
speed of the reference point is
\(u_1 = v + d_\mathrm{ccw}\,\omega^+ + d_\mathrm{cw}\,\omega^-\) with
\(\omega^\pm = \max(\pm\omega, 0)\). Since
\(|\omega| = \omega^+ + \omega^-\) and \(\omega = \omega^+ - \omega^-\),
\begin{equation}
d_\mathrm{ccw}\,\omega^+ + d_\mathrm{cw}\,\omega^- = \bar d\,|\omega| + \tfrac12 (d_\mathrm{ccw} - d_\mathrm{cw})\,\omega,
\label{eq:cor}
\end{equation}
with \(\bar d = \tfrac12 (d_\mathrm{ccw} + d_\mathrm{cw})\).
In the forward map \(\omega\) is linear in \(\boldsymbol\phi\), and so
is \(v\). The antisymmetric term is therefore absorbed exactly by
\(v(\boldsymbol\phi)\): any pair \((d_\mathrm{ccw}, d_\mathrm{cw})\)
with the same mean fits steady-state data equally well, with a
correspondingly shifted \(v\). Only \(\bar d\) is identifiable, through
the kink of \(|\omega|\) at \(\omega = 0\), which requires both turning
directions in the grid. The \texttt{cor} variant fits
\(u_1 = v(\boldsymbol\phi) + \bar d\,|\hat\omega|\) with \(\hat\omega\)
the model's own yaw rate and exports \(\bar d\); on a simulated robot
with the twin's 14.5 / 16.1 mm split it recovers 15.3 mm, the mean. The
design matrix \([1, f_L, f_R, \omega^+, \omega^-]\) built with the
model's own \(\hat\omega\) has rank 4 of 5 on the lab's tables (with
the measured \(\omega\) it has full rank only through the noise). The twin's split is therefore a modelling choice that
steady-state data cannot confirm; whether transient data break the
degeneracy is an open question we test in §\ref{sec:protocol} \hw{step-response fit on
robot \#1: are separate CW and CCW offsets identified, with CIs}. More
generally we report the Fisher information of the fitted model and the
parameter standard deviations against the Cramér--Rao bound, as Censi et
al. \cite{censi2013calibration} do for odometry, and the condition
number of the design, which exposes the trade-off between \(p_0\) and
\(\gamma\) on grids that start far above the dead zone.

\subsection{Inverse map and feasible
set}\label{sec:inverse}

For the \texttt{lin} and \texttt{cor} variants \((v, \omega)\) is affine
in \((f_L, f_R)\), so a commanded \((v,\omega)\) is inverted by a 2×2
solve followed by the inverse activation
\(p_k = p_{0,k} + f_k^{1/\gamma_k}\), clipped to the drive range. The
image of the activation box under the affine part is a parallelogram
(four constraints \(a_i v + b_i \omega \le c_i\)) that we export as the
feasible set; targets outside it are projected. The whole inverse
fits in a few hundred bytes of C and needs no lookup table.

\subsection{Lag, dead time and noise}\label{sec:lag}

Steps and pseudo-random binary sequences switching every 0.1--0.5 s, cut
into chunks between re-centrings, give a first-order lag \(\tau\) and
dead time per robot, by fitting the integrated response as in Seegmiller
et al. \cite{seegmiller2013ipem}. Sixty-second clockwise and
counter-clockwise circles give the heading diffusion \(D_r\) (Brownian
or Ornstein--Uhlenbeck, whichever the autocorrelation supports), speed
jitter and lateral slip; the mean-squared displacement is fitted on a
log scale with bootstrap intervals \cite{bailey2022logit}. A 60 s
static recording gives the tracker's noise floor.

\subsection{Population model, prior and few-shot
fit}\label{sec:prior}

Let
\(\boldsymbol\theta = (p_{0,L}, p_{0,R}, \gamma_L, \gamma_R, \mathrm{vec}\,\mathbf C)\).
From \(S\) fully calibrated robots we form the population prior, using
each robot's day-averaged \(\boldsymbol\theta_r\) (the day level
\(\boldsymbol\Sigma_d\) is estimated separately in
§\ref{sec:variance}),
\(\boldsymbol\theta \sim \mathcal N(\boldsymbol\mu, \boldsymbol\Sigma)\)
with \(\boldsymbol\mu\) the sample mean and \(\boldsymbol\Sigma\) the
sample covariance shrunk half-way towards its diagonal and floored by a
diagonal term that shrinks as \(1/\sqrt S\); for \(S = 1\) the prior is
the reference robot's parameters with the floor alone, the ``one
reference robot with inflated variance'' fallback. Cross-covariance
between the nonlinear block \((p_0, \gamma)\) and the coefficient rows,
and between rows, is dropped because it is not estimable from a handful
of robots; the price is a prior that is less informative and may
misstate its confidence along correlated directions. The segment noise
\(\sigma_o\) per output is the residual of the full fits.

A new robot with \(k\) segments is fitted by maximising the posterior,
\begin{multline}
\hat{\boldsymbol\theta} = \arg\min_{\boldsymbol\theta}\ \sum_{i=1}^{k} \sum_{o} \frac{\bigl(y_{io} - F_o(L_i,R_i,V_i;\boldsymbol\theta)\bigr)^2}{\sigma_o^2} \\
+ (\boldsymbol\theta-\boldsymbol\mu)^\top \boldsymbol\Sigma^{-1} (\boldsymbol\theta-\boldsymbol\mu),
\label{eq:map}
\end{multline}
by the same variable projection as §\ref{sec:forward-map}: given \((p_0,\gamma)\) each
coefficient row is a Gaussian-prior least-squares problem with a
closed-form solution and posterior covariance, and Nelder--Mead searches
the four nonlinear parameters from the prior mean. With \(k = 0\) the
estimate is the prior mean; with \(k\) large it approaches the full fit.
The posterior covariance of the rows gives a predictive interval for the
mean twist at every drive cell, whose coverage we report.

\subsection{Which segments to
measure}\label{sec:design}

Given the prior, the information a segment at cell \(c\) adds is
\(\mathbf J_c^\top \mathbf R^{-1} \mathbf J_c\), with \(\mathbf J_c\)
the \(3 \times \dim\boldsymbol\theta\) Jacobian of the twist at the
prior mean and \(\mathbf R = \mathrm{diag}(\sigma_o^2)\). We choose
segments greedily to maximise
\(\log\det(\boldsymbol\Sigma^{-1} + \sum_c \mathbf J_c^\top \mathbf R^{-1} \mathbf J_c)\)
(Bayesian D-optimality
\cite{chaloner1995design, sun2008observability}), allowing repeated
cells. Because the model is nonlinear the design is local to the prior
mean; in the simulation of §\ref{sec:simulation} the chosen cells are the single-motor
pivots, full power on both motors, and cells at the prior's dead-zone
edge. The alternative of refitting after each segment and
re-designing \cite{rainforth2024design} is a straightforward extension.

\subsection{Variance decomposition}\label{sec:variance}

With every robot measured on \(D\) days with \(n\) replicates per cell,
each cell and output is a balanced nested design (robot ⊃ day ⊃
segment). Method-of-moments estimates of the three variance components,
averaged over the cells with motion, give the robot, day and segment
standard deviations per output and their fractions of the total. The
robot component is what the population prior needs; the day component is
drift; the segment component is the noise floor of one measurement.

\subsection{Baselines}\label{sec:baselines}

For the few-shot study we compare (i) \texttt{hier}, the MAP fit above
with a prior from \(S \in \{1, 2, 4\}\) robots; (ii)
\texttt{scratch}, the full variable-projection fit on the \(k\) segments
alone, which is undefined below about eight segments; (iii)
\texttt{ref}, the prior mean used without data; (iv) \texttt{ref\_gain},
the prior mean corrected by a per-output affine map
\(y_o = a_o\,\hat y_o + b_o\) fitted on the \(k\) segments, the
master/slave standardisation of §\ref{sec:rw-transfer}; and (v) a Klingner-style
black-box search \cite{klingner2014droplet} over the same segment
budget \hw{run on hardware only; not in the simulation}. Segments are chosen at random from the robot's
measured pool or by the D-optimal design, and every held-out robot is
scored on its measured segments outside the chosen subset (in
simulation, on the noise-free truth over the whole grid).
